\section{Livrable 1 et 2 -- Synoptique et architecture réseau}

\begin{UPSTIactivite}[][Schéma synoptique de l'installation]
    \UPSTIquestion{À partir de la présentation du système, identifier
    l'ensemble des équipements physiques constituant la cellule robotique
    (automate, robot, convoyeur, IHM, cuves, capteurs et actionneurs
    nécessaires, etc.).}

    \UPSTIquestion{Proposer un \textbf{schéma synoptique} de l'installation
    faisant apparaître ces équipements ainsi que les liaisons fonctionnelles
    entre eux (échanges d'informations, flux de pièces).}
\end{UPSTIactivite}

\begin{UPSTIactivite}[][Architecture réseau et plan d'adressage]
    \UPSTIquestion{En tenant compte de la contrainte de séparation réseau
    (\cref{secContraintes}), proposer une \textbf{architecture réseau}
    permettant de relier les équipements nécessitant une communication
    EtherNet/IP entre eux.}

    \UPSTIquestion{Identifier les équipements qui doivent obligatoirement se
    trouver sur le réseau séparé, et ceux qui peuvent rester sur le réseau de
    l'entreprise.}

    \UPSTIquestion{Proposer un \textbf{plan d'adressage IP} cohérent pour le
    réseau séparé, en justifiant le choix du plan d'adresses et du masque de
    sous-réseau retenus. Ce plan devra rester cohérent avec l'adressage déjà
    utilisé par les ordinateurs de la salle.}

    \UPSTIquestion{Représenter un \textbf{schéma de l'architecture réseau}
    complet, faisant apparaître pour chaque équipement son adresse IP et son
    rôle (scanner, adapter, etc. au sens EtherNet/IP -- voir le document
    ressource EtherNet/IP).}
\end{UPSTIactivite}
